% Worked example, alignment and prosody marking: one panel on the time axis of the clip
% (1 s = 1.45 cm), with one column per word. Each column holds the word's aligned block (below
% the audio envelope), its name, its loudness bar and its length bar, joined by a faint vertical
% guide. Names are vertical (90 degrees) and small, hanging below their blocks. Row titles are on
% the left and the legend is on the right. Row positions (y, cm):
%   audio envelope 0, word blocks -0.55, names from -0.80 down,
%   loudness baseline -3.00, length baseline -5.70, time axis -5.85.
\begin{tikzpicture}[x=1.45cm,
    emph none/.style={fill=black!25},
    emph loud/.style={fill=red!60},
    emph soft/.style={fill=blue!50},
    drag no/.style={fill=black!25},
    drag yes/.style={fill=orange!75},
    flag none/.style={draw=black!50},
    flag gap/.style={draw=red!75!black, line width=1pt},
    flag short/.style={flag gap},
    label none/.style={},
    label gap/.style={text=red!75!black},
    label short/.style={label gap},
    panel/.style={fill=black!3, draw=black!15, rounded corners=3pt},
    row title/.style={font=\footnotesize\bfseries, align=right, anchor=east},
    legend/.style={font=\scriptsize, anchor=west},
    guide/.style={black!22, densely dotted},
  ]
  \pgfmathtruncatemacro\lastsecond{\clipduration}
  \def\legendx{6.75}   % where the legend starts (s)
  \def\wordy{-0.55}    % centre of the word blocks
  \def\loudy{-3.00}    % loudness baseline (clip average)
  \def\lengthy{-5.70}  % length baseline (zero length)

  % One panel for all rows
  \fill[panel] (-1.55,0.55) rectangle (8.95,-6.45);

  % Row titles
  \node[row title] at (-0.1,-0.55)   {Words and\\pauses};
  \node[row title] at (-0.1,\loudy)  {Loudness};
  \node[row title] at (-0.1,-5.05)   {Length vs.\\expected};

  % Faint row separators
  \draw[black!12] (-1.4,-2.15) -- (8.8,-2.15);
  \draw[black!12] (-1.4,-4.10) -- (8.8,-4.10);

  % Guides: one per word, from its block down through the loudness and length rows
  \foreach \word/\tstart/\tend/\flag/\loudz/\emphasis/\ratio/\dragged in \words
    \draw[guide] ({(\tstart+\tend)/2},{\wordy-0.18}) -- ({(\tstart+\tend)/2},\lengthy);

  % Reference lines
  \draw[black!40] (0,\loudy) -- (\clipduration,\loudy);                       % clip average loudness
  \draw[black!40] (0,\lengthy) -- (\clipduration,\lengthy);                   % zero length
  \draw[black!50, dashed] (0,{\lengthy+0.6}) -- (\clipduration,{\lengthy+0.6}); % expected length

  % Audio envelope of the clip, above the word blocks
  \begin{scope}[yscale=0.3]
    \fill[black!20] plot file {figures/data/original-envelope.dat} -- cycle;
  \end{scope}

  % One column per word; words that follow a silence of 0.5 s or more are outlined in red
  \foreach \word/\tstart/\tend/\flag/\loudz/\emphasis/\ratio/\dragged in \words {
    \draw[flag \flag, fill=black!8] (\tstart,{\wordy-0.18}) rectangle (\tend,{\wordy+0.18});
    % the vertical name hangs below the block and is filled with the panel colour,
    % so that it masks the guide behind it
    \node[font=\tiny\ttfamily, rotate=90, anchor=east, inner sep=0.8pt, fill=black!3, label \flag]
      at ({(\tstart+\tend)/2},{\wordy-0.22}) {\word};
    % loudness: bar height follows the z-score (0.3 cm per standard deviation)
    \fill[emph \emphasis] (\tstart,\loudy) rectangle (\tend,{\loudy + 0.3*\loudz});
    % length: 0.6 cm corresponds to the expected length
    \fill[drag \dragged] (\tstart,\lengthy) rectangle (\tend,{\lengthy + 0.6*\ratio});
  }

  % Pause markers: a red bar with its length for a long pause, an orange comma for a short one
  \def\pauselong#1#2{%
    \draw[red!75!black, line width=1.6pt] (#1,{\wordy-0.32}) -- (#1,{\wordy+0.32});
    \node[font=\scriptsize, text=red!75!black, below] at (#1,{\wordy-0.32}) {#2\,s};}
  \def\pauseshort#1#2{%
    \node[font=\Large\bfseries, text=orange!80!black] at (#1,{\wordy-0.12}) {,};}
  \foreach \at/\type/\gap in \pauses {
    \csname pause\type\endcsname{\at}{\gap}
  }

  % Label the word with the largest length ratio
  \gdef\maxratio{0}
  \foreach \word/\tstart/\tend/\flag/\loudz/\emphasis/\ratio/\dragged in \words {
    \pgfmathparse{\ratio>\maxratio ? 1 : 0}
    \ifnum\pgfmathresult=1
      \xdef\maxratio{\ratio}
      \pgfmathsetmacro\middle{(\tstart+\tend)/2}
      \xdef\maxat{\middle}
    \fi
  }
  \node[font=\scriptsize, text=orange!60!black, above] at (\maxat,{\lengthy + 0.6*\maxratio})
    {\pgfmathprintnumber[fixed, precision=1]{\maxratio}×};

  % Time axis, below the length row
  \draw[black!60] (0,-5.85) -- (\clipduration,-5.85);
  \foreach \t in {0,...,\lastsecond}
    \draw[black!60] (\t,-5.85) -- ++(0,-0.08) node[below, font=\scriptsize] {\t\,s};

  % Legend, grouped beside each row
  \begin{scope}[shift={(\legendx,0)}, x=1cm]
    % words and pauses
    \fill[black!20] (0,0.20) rectangle (0.25,0.40);                 \node[legend] at (0.3,0.30) {audio envelope};
    \draw[black!50, fill=black!8] (0,-0.20) rectangle (0.25,0.00); \node[legend] at (0.3,-0.10) {word};
    \draw[red!75!black, line width=1pt, fill=black!8] (0,-0.70) rectangle (0.25,-0.50);
    \node[legend, align=left] at (0.3,-0.60) {word after a silence\\of 0.5\,s or more};
    \draw[red!75!black, line width=1.6pt] (0.12,-1.30) -- (0.12,-1.00); \node[legend] at (0.3,-1.15) {long pause (length)};
    \node[font=\Large\bfseries, text=orange!80!black] at (0.12,-1.65) {,}; \node[legend] at (0.3,-1.60) {short pause};
    % loudness
    \fill[red!60]   (0,-2.60) rectangle (0.25,-2.40); \node[legend] at (0.3,-2.50) {loud};
    \fill[blue!50]  (0,-2.95) rectangle (0.25,-2.75); \node[legend] at (0.3,-2.85) {soft};
    \fill[black!25] (0,-3.30) rectangle (0.25,-3.10); \node[legend] at (0.3,-3.20) {unmarked};
    \draw[black!40] (0,-3.55) -- (0.25,-3.55);        \node[legend] at (0.3,-3.55) {clip average};
    % length
    \fill[orange!75] (0,-4.50) rectangle (0.25,-4.30); \node[legend] at (0.3,-4.40) {dragged};
    \fill[black!25]  (0,-4.85) rectangle (0.25,-4.65); \node[legend] at (0.3,-4.75) {not dragged};
    \draw[black!50, dashed] (0,-5.10) -- (0.25,-5.10); \node[legend] at (0.3,-5.10) {expected length};
  \end{scope}
\end{tikzpicture}
